\section{Discussion}
The results suggest that StratEval is useful not because it assigns a single global safety score, but because it makes strategic failure modes more legible. The evaluation compares recognition-oriented probes with decision-stage behavior under matched scenario frames, which helps reveal when a model is merely confused and when its decision response remains strategically misaligned despite evidence of scenario awareness.

\paragraph{Capability and strategic failure}
One of the clearest patterns in the results is a divergence between situational awareness and behavioral restraint. \texttt{Gemma-3-12B} was better at identifying the operative conflict in a scenario, but that higher recognition rate did not reduce severe escalation. In many cases, its severe decision responses were also more coherent and legible. This reinforces a broader concern in alignment research: better reasoning and planning do not automatically imply better alignment, and may instead make strategic failure easier to formulate and harder to dismiss as mere confusion.

\paragraph{Surface alignment versus strategic grounding}
The decomposed prompting protocol also highlights the distinction between surface alignment and deeper behavioral stability. Under the same scenario frame, a model may correctly explain why blackmail is unethical, or accurately summarize that a shutdown order conflicts with its task, while its decision-stage response still endorses extortion, deception, or procedural sabotage. In our stateless results, this is a matched-prompt comparison rather than a remembered sequence inside one continuous conversation. That distinction matters, but it does not dissolve the concern: in this two-model comparison, the model with the higher situational-awareness rate did not have a lower severe-escalation rate in the decision-stage outputs, and many severe decision responses were too coherent to dismiss as simple misunderstanding. This is precisely why one-shot safety tests are insufficient for the class of failures considered here. An evaluation that only asks whether the model can state the correct principle may overestimate alignment if the model's abstract recognition is weakly coupled to its decision behavior.

What StratEval makes visible is the gap between \emph{moral description} and \emph{behavioral commitment}. Recognition-stage success shows that the model can represent the conflict, identify the prohibited tactic, and sometimes even articulate why that tactic is wrong. But, in the matched decision condition, that recognition evidence does not guarantee restraint. In that sense, the failure is not simple ignorance. It is a form of strategic instability in which the model's explicit account of the norms governing the situation is not reliably coupled to its chosen course of action. The issue, then, is not merely whether a model can state the rule, but whether it remains behaviorally constrained when the scenario rewards breaking it. Stronger claims about a model carrying its own recognition forward and then reasoning around it require the hybrid or stateful modes identified as future work.

\paragraph{Why taxonomy matters}
The taxonomic framework is valuable because it prevents all failures from collapsing into a single undifferentiated category. A polite but instrumentally evasive response, a deceptive concealment strategy, a direct coercive threat, and a multi-step scheme to subvert oversight are all distinct kinds of failure. They likely call for different forms of intervention, different forms of monitoring, and different deployment restrictions. StratEval shows that generalized statements and predictions concerning AI safety may be inadequate, and importantly shifts the discussion to something more concrete: what might models actually do when placed under instrumental convergence pressure?

\paragraph{Implications for evaluation design}
A further lesson is methodological. Alignment evaluations should not assume that final-answer correctness is enough. For strategic safety questions, it is often important to know whether the model can recognize the relevant conflict, understand the tactic in question, and articulate the constraints around it under the same scenario frame used for decision-stage evaluation. That background does not prove intent in any strong philosophical sense, but it makes severe misaligned decisions far more interpretable.

\paragraph{Limitations}
This study has several important limitations. First, the reported headline results are restricted to two open-weight model families and to the stateless modality. The present paper therefore should not be read as a broad census of model behavior across architectures, training regimes, or deployment settings. The results are best understood as evidence that strategically concerning behavior can be elicited clearly and repeatedly in this evaluation frame, not as a final estimate of prevalence across all contemporary systems.

The stateless design also limits the interpretation of the recognition-to-decision comparison. Because each probe re-presents the full scenario and does not carry the model's previous answer forward, the headline results do not by themselves show that a single remembered chain of thought first recognized the conflict and then chose to violate it. They show instead that, under matched scenario frames, recognition-stage competence can coexist with severe decision-stage escalation. Hybrid and stateful evaluations are required to test stronger within-trajectory claims.

Second, the scenarios are narrative and text-mediated. They probe strategic reasoning, role-conditioned planning, and whether outputs endorse harmful or self-protective action in language, but they do not directly test embodied capability or real-world execution. A model that describes an instrumental strategy is not identical to an agent that can robustly carry it out in a live environment. At the same time, we do not regard this as a trivial limitation: for alignment work, the production of such plans under pressure is itself an important warning signal, even when the surrounding environment is only simulated textually.

Third, the taxonomy and escalation ladder are rubric-guided judgments produced by an LLM-assisted adjudication pipeline with human spot-checks, not direct ground truth about internal intent. This means that the results depend partly on the quality of the judge, the clarity of the rubric, and the stability of the taxonomy. The judge-variance findings reported earlier make this especially important: weaker judges can under-read strategically grounded behavior. The evaluation therefore does not eliminate interpretive uncertainty; it organizes and disciplines it.

Fourth, although the scenarios vary along multiple pressure-related axes, the present paper does not claim a fully validated scalar model of ``pressure'' as such. The evaluation was constructed to vary shutdown threat, leverage, stakes, access, concealment opportunity, and related factors, and the resulting corpus does span a meaningful gradient of strategic stress. But the paper does not establish that these factors collapse cleanly into a single psychometrically validated latent variable. We therefore treat ``pressure'' here as a design concept and descriptive framing device, not a finished measurement theory.

Finally, because the evaluation is explicitly constructed to surface dangerous strategic behavior, it remains possible that some of the prompting structure itself helps make those behaviors salient. We regard this as an acceptable tradeoff for diagnostic evaluation, but it does place a boundary on interpretation. StratEval is not designed to model every ordinary deployment context. It is designed to stress the regime in which oversight, incentives, and strategic opportunity come into conflict, and to make the resulting failure modes legible.

\paragraph{Future work}
Several extensions are natural. The most immediate is to complete broader hybrid and stateful sweeps so that context accumulation can be studied systematically. Another is to expand the scenario corpus into additional operational domains, especially ones involving tool use, code execution, or richer human-in-the-loop interaction. More broadly, future work should test whether targeted interventions can reduce specific failure clusters---for example, deception-heavy outputs, procedural capture, or self-preservation reasoning---rather than treating alignment as a monolithic property.

In short, StratEval does not primarily show that a model is globally ``good'' or ``bad.'' It shows, more usefully, how strategic failure can manifest, how severe those failures can become, and why alignment evaluation benefits from structured failure analysis rather than summary judgment alone.
